\section*{Capítulo \ref{ch:g7:identifying-substances} --- Identificar substâncias}

\begin{solution}{exo:g7:identifying-substances:1}
O teste da água de cal: o gás é borbulhado numa água de cal límpida, que
fica leitosa se o gás for dióxido de carbono.
\end{solution}

\begin{solution}{exo:g7:identifying-substances:2}
O pó é sulfato de cobre anidro, e o líquido contém água.
\end{solution}

\begin{solution}{exo:g7:identifying-substances:3}
O oxigênio.
\end{solution}

\begin{solution}{exo:g7:identifying-substances:4}
O grafite do lápis não \omterm{def:g2:mixing-and-dissolving:dissolve}{se dissolve} no \omterm{def:g6:solutions-and-solubility:solute}{solvente}. Uma linha de tinta seria
ela mesma levada para cima e separada, e estragaria o \omterm{def:g7:identifying-substances:chromatography}{cromatograma}.
\end{solution}

\begin{solution}{exo:g7:identifying-substances:5}
Se a mancha ficasse abaixo do \omterm{def:g6:solutions-and-solubility:solute}{solvente}, os corantes se dissolveriam no
\omterm{def:g6:solutions-and-solubility:solute}{solvente} do béquer em vez de subir pelo papel.
\end{solution}

\begin{solution}{exo:g7:identifying-substances:6}
Não, ela é uma \omterm{def:g6:pure-substances-and-mixtures:mixture}{mistura} de dois corantes: provavelmente o corante amarelo
e o corante azul das referências.
\end{solution}

\begin{solution}{exo:g7:identifying-substances:7}
Recolher um pouco do gás num tubo de ensaio e segurar uma tala acesa na
boca do tubo: um estalo agudo mostra que é hidrogênio.
\end{solution}

\begin{solution}{exo:g7:identifying-substances:8}
O \omterm{def:g4:air-a-mixture-of-gases:air}{ar} expirado contém muito mais dióxido de carbono do que o \omterm{def:g4:air-a-mixture-of-gases:air}{ar} fresco.
\end{solution}

\begin{solution}{exo:g7:identifying-substances:9}
O teste mostra que há água, não que ela está sozinha: água salgada, suco
ou qualquer \omterm{def:g6:pure-substances-and-mixtures:mixture}{mistura} aquosa também deixariam o pó azul.
\end{solution}

\begin{solution}{exo:g7:identifying-substances:10}
Três corantes (três manchas). O corante amarelo foi o que subiu mais alto.
\end{solution}

\begin{solution}{exo:g7:identifying-substances:11}
Por exemplo: uma tala em brasa em cada frasco --- o que a reacende
contém oxigênio. Nos outros dois, uma tala acesa na boca --- um estalo
agudo mostra o hidrogênio. O último frasco contém dióxido de carbono (o
que pode ser confirmado com água de cal). Dois ou três testes bastam.
\end{solution}

\begin{solution}{exo:g7:identifying-substances:12}
O gás pode não conter dióxido de carbono, ou conter pouco demais para
aparecer nesse tempo. Um teste com \omterm{def:g4:air-a-mixture-of-gases:air}{ar} expirado, que sabidamente contém
dióxido de carbono, verifica se a água de cal está funcionando.
\end{solution}

\begin{solution}{pb:g7:identifying-substances:1}
\textbf{1.} Para que as condições (papel, \omterm{def:g6:solutions-and-solubility:solute}{solvente}, tempo) sejam as
mesmas para as quatro: só assim as alturas das manchas podem ser
comparadas.

\textbf{2.} Uma \omterm{def:g6:pure-substances-and-mixtures:mixture}{mistura}: a tinta dele dá três manchas separadas, três
corantes.

\textbf{3.} Dois corantes. Não: a tinta do bilhete tem três corantes, e
a caneta 1 não tem nem o corante vermelho nem o amarelo.

\textbf{4.} As três manchas dela não estão nas mesmas alturas que as do
bilhete: são três corantes diferentes.

\textbf{5.} A caneta 2: as três manchas dela estão exatamente nas
alturas das três manchas do bilhete.

\textbf{6.} Oxigênio.

\textbf{7.} Hidrogênio.

\textbf{8.} Dióxido de carbono.

\textbf{9.} Um branco (um teste de controle). Ele mostra que a água de
cal não fica leitosa sozinha, nem com \omterm{def:g4:air-a-mixture-of-gases:air}{ar} comum nesse tempo: logo, a
turvação com o gás Z vem de fato do dióxido de carbono.

\textbf{10.} ``Água --- sulfato de cobre anidro --- o branco fica
azul'': o líquido contém água.

\textbf{11.} Não. Um teste positivo mostra que há água, não que não há
mais nada: o líquido poderia ser uma solução ou outra \omterm{def:g6:pure-substances-and-mixtures:mixture}{mistura} aquosa.

\textbf{12.} A tinta do bilhete contém 3 corantes, e foi a caneta 2 que
escreveu o bilhete.
\end{solution}
